% Generated draft: requires mathematical and source review.
\documentclass[11pt]{article}
\usepackage[margin=25mm]{geometry}
\usepackage{fontspec}
\setmainfont{Latin Modern Roman}
\usepackage{amsmath,amssymb,mathtools,unicode-math}
\setmathfont{Latin Modern Math}
\usepackage{enumitem,needspace,xurl,hyperref}
\setlength{\parindent}{0pt}
\setlength{\parskip}{5pt}
\setlength{\emergencystretch}{3em}
\newcommand{\N}{\mathbb{N}}
\newcommand{\Z}{\mathbb{Z}}
\newcommand{\Q}{\mathbb{Q}}
\newcommand{\R}{\mathbb{R}}
\newcommand{\C}{\mathbb{C}}
\newcommand{\F}{\mathbb{F}}
\newcommand{\A}{\mathbb{A}}
\newcommand{\PP}{\mathbb{P}}
\newcommand{\E}{\mathbb{E}}
\newcommand{\eps}{\varepsilon}
\newcommand{\dd}{\,\mathrm{d}}
\begin{document}
\section*{UnsolvedMath ID: 701434}
\textbf{Source problem code:} catalogue entry
\subsection{Definitions and mathematical statement}
A group \(H\) is soluble if its derived series, defined by \(H^{(0)}=H\) and \(H^{(i+1)}=[H^{(i)},H^{(i)}]\), reaches the trivial group after finitely many steps. For a soluble group with \(H^{(d)}=1\), interpret its torsion-free rank as its Hirsch length
\[
r_0(H)=\sum_{i=0}^{d-1}\dim_{\Q}\!\left(\Q\otimes_{\Z}\bigl(H^{(i)}/H^{(i+1)}\bigr)\right).
\]
The source does not explicitly define \(r_0\); this is the standard soluble-group interpretation used in this draft.

\textbf{Problem.} For every soluble group \(G\) with \(r_0(G)<\infty\), and every pair of subgroups \(A,B\leq G\) satisfying the setwise factorization
\[
G=AB=\{ab:a\in A,\ b\in B\},
\]
is it true that
\[
r_0(G)=r_0(A)+r_0(B)-r_0(A\cap B)?
\]
No normality or finite generation of the subgroups is assumed. The source states that the inequality
\[
r_0(G)\leq r_0(A)+r_0(B)-r_0(A\cap B)
\]
always holds in the stated setting; this is a supplied background assertion.

\textbf{Unresolved qualification.} The original wording is ``the (soluble) group''. This draft explicitly states the soluble version. Whether the parentheses indicate an additional intended question for groups that are not soluble is not determined by the supplied record.
\subsection{Short English statement}
For a soluble group of finite torsion-free rank that is the setwise product of two subgroups, must its rank equal the sum of their ranks minus the rank of their intersection? The source supplies the upper-bound inequality. The parenthetical solubility qualification and the intended definition of rank require confirmation.
\subsection{Sources}
\begin{itemize}
\item Supplied local UnsolvedMath registry record, permanent ID \(701434\), exact source problem code \texttt{catalogue entry}, title \emph{Kourovka Notebook Problem 18.5}. The registry labels the problem \texttt{open}; this is supplied metadata, not an independently verified mathematical assessment.
\item B. Amberg, proposed in 2014, Problem 18.5 in \emph{The Kourovka Notebook (Unsolved Problems in Group Theory)}, 18th issue (2014). The supplied background also identifies the 21st issue edition (2026), \texttt{arXiv:1401.0300v48}. The supplied collection description names E. I. Khukhro and V. D. Mazurov as editors. These bibliographic details are supplied attribution and have not been independently checked.
\item Registry record updated \texttt{2026-10-16T14:00:00Z}; upstream revision \texttt{d371ee917ac9647fbb06cd6fc529deec6f3dda82}. This timestamp is a record update, not a supplied literature-verification date.
\end{itemize}
\end{document}
